%% errors.tex
%%
%% Copyright (C) 2000-2018 Simone Piccardi.  Permission is granted to
%% copy, distribute and/or modify this document under the terms of the GNU Free
%% Documentation License, Version 1.1 or any later version published by the
%% Free Software Foundation; with the Invariant Sections being "Un preambolo",
%% with no Front-Cover Texts, and with no Back-Cover Texts.  A copy of the
%% license is included in the section entitled "GNU Free Documentation
%% License".
%%

\chapter{I codici di errore}
\label{cha:errors}

Si riportano in questa appendice tutti i codici di errore. Essi sono
accessibili attraverso l'inclusione del file di header \headfile{errno.h}, che
definisce anche la variabile globale \var{errno}. Per ogni errore definito
riporteremo la stringa stampata da \func{perror} ed una breve spiegazione. Si
tenga presente che spiegazioni più particolareggiate del significato
dell'errore, qualora necessarie per casi specifici, possono essere trovate
nella descrizione del prototipo della funzione per cui detto errore si è
verificato.

I codici di errore sono riportati come costanti di tipo \ctyp{int}, i valori
delle costanti sono definiti da macro di preprocessore nel file citato, e
possono variare da architettura a architettura; è pertanto necessario
riferirsi ad essi tramite i nomi simbolici. Le funzioni \func{perror} e
\func{strerror} (vedi sez.~\ref{sec:sys_strerror}) possono essere usate per
ottenere dei messaggi di errore più espliciti.


\section{Gli errori dei file}
\label{sec:err_file_errors}

In questa sezione sono raccolti i codici restituiti dalle \textit{system call}
attinenti ad errori che riguardano operazioni specifiche relative alla
gestione dei file.

\begin{basedescript}{\desclabelwidth{1.5cm}\desclabelstyle{\nextlinelabel}}
\item[\errcode{EACCES} \textit{Permission denied}.] Permesso negato; l'accesso
  al file o alla directory non è consentito: i permessi del file o della
  directory o quelli necessari ad attraversare un \textit{pathname} non
  consentono l'operazione richiesta. 
\item[\errcode{EBADF} \textit{Bad file descriptor}.] File descriptor non
  valido: si è usato un file descriptor inesistente, o aperto in sola lettura
  per scrivere, o viceversa, o si è cercato di eseguire un'operazione non
  consentita per quel tipo di file descriptor.
\item[\errcode{EBUSY} \textit{Resource busy}.] Una risorsa di sistema che non
  può essere condivisa è occupata. Ad esempio si è tentato di cancellare la
  directory su cui si è montato un filesystem.
\item[\errcode{EDQUOT} \textit{Quota exceeded}.] Si è ecceduta la quota di
  disco dell'utente, usato sia per lo spazio disco che per il numero di
  \textit{inode}.
\item[\errcode{EEXIST} \textit{File exists}.] Si è specificato un file
  esistente in un contesto in cui ha senso solo specificare un nuovo file.
\item[\errcode{EFBIG} \textit{File too big}.] Si è ecceduto il limite imposto
  dal sistema sulla dimensione massima che un file può avere.
\item[\errcode{EFTYPE} \textit{Inappropriate file type or format}.] Il file è
  di tipo sbagliato rispetto all'operazione richiesta o un file di dati ha un
  formato sbagliato. Alcuni sistemi restituiscono questo errore quando si
  cerca di impostare lo \textit{sticky bit} su un file che non è una
  directory.
\item[\errcode{EIO} \textit{Input/output error}.] Errore di input/output:
  usato per riportare errori hardware in lettura/scrittura su un dispositivo.
\item[\errcode{EISDIR} \textit{Is a directory}.] Il file specificato è una
  directory; non può essere aperto in scrittura, né si possono creare o
  rimuovere link diretti ad essa.
\item[\errcode{ELOOP} \textit{Too many symbolic links encountered}.] Ci sono
  troppi link simbolici nella risoluzione di un \textit{pathname}.
\item[\errcode{EMFILE} \textit{Too many open files}.] Il processo corrente ha
  troppi file aperti e non può aprirne altri. Anche i descrittori duplicati ed
  i socket vengono tenuti in conto.\footnote{il numero massimo di file aperti
    è controllabile dal sistema; in Linux si può impostare usando il comando
    \cmd{ulimit}, esso è in genere indicato dalla costante \const{OPEN\_MAX},
    vedi sez.~\ref{sec:sys_limits}.}
\item[\errcode{EMLINK} \textit{Too many links}.] Ci sono già troppi link al
  file, il numero massimo è specificato dalla variabile \const{LINK\_MAX},
  vedi sez.~\ref{sec:sys_limits}.
\item[\errcode{ENAMETOOLONG} \textit{File name too long}.] Si è indicato un
  \textit{pathname} troppo lungo per un file o una directory.
\item[\errcode{ENFILE} \textit{File table overflow}.] Il sistema ha troppi
  file aperti in contemporanea. Si tenga presente che anche i socket contano
  come file. Questa è una condizione temporanea, ed è molto difficile che si
  verifichi nei sistemi moderni.
\item[\errcode{ENODEV} \textit{No such device}.] Si è indicato un tipo di
  device sbagliato ad una funzione che ne richiede uno specifico.
\item[\errcode{ENOENT} \textit{No such file or directory}.] Il file indicato
  dal \textit{pathname} non esiste: o una delle componenti non esiste o il
  \textit{pathname} contiene un link simbolico spezzato.  Errore tipico di un
  riferimento ad un file che si suppone erroneamente essere esistente.
\item[\errcode{ENOEXEC} \textit{Invalid executable file format}.] Il file non
  ha un formato eseguibile, è un errore riscontrato dalle funzioni
  \func{exec}.
\item[\errcode{ENOLCK} \textit{No locks available}.] È usato dalle utilità per
  la gestione del file locking; non viene generato da un sistema GNU, ma può
  risultare da un'operazione su un server NFS di un altro sistema.
\item[\errcode{ENOSPC} \textit{No space left on device}.] La directory in cui
  si vuole creare il link non ha spazio per ulteriori voci, o si è cercato di
  scrivere o di creare un nuovo file su un dispositivo che è già pieno.
\item[\errcode{ENOTBLK} \textit{Block device required}.] Si è specificato un
  file che non è un \textit{block device} in un contesto in cui era necessario
  specificare un \textit{block device} (ad esempio si è tentato di montare un
  file ordinario).
\item[\errcode{ENOTDIR} \textit{Not a directory}.] Si è specificato un file
  che non è una directory in una operazione che richiede una directory.
\item[\errcode{ENOTEMPTY} \textit{Directory not empty}.] La directory non è
  vuota quando l'operazione richiede che lo sia. È l'errore tipico che si ha
  quando si cerca di cancellare una directory contenente dei file.
\item[\errcode{ENOTTY} \textit{Not a terminal}.] Si è tentata una operazione
  di controllo relativa ad un terminale su un file che non lo è.
\item[\errcode{ENXIO} \textit{No such device or address}.] Dispositivo
  inesistente: il sistema ha tentato di usare un dispositivo attraverso il
  file specificato, ma non lo ha trovato. Può significare che il file di
  dispositivo non è corretto, che il modulo relativo non è stato caricato nel
  kernel, o che il dispositivo è fisicamente assente o non funzionante.
\item[\errcode{EPERM} \textit{Operation not permitted}.] L'operazione non è
  permessa: solo il proprietario del file o un processo con sufficienti
  privilegi può eseguire l'operazione.
\item[\errcode{EPIPE} \textit{Broken pipe}.] Non c'è un processo che stia
  leggendo l'altro capo della \textit{pipe}. Ogni funzione che restituisce
  questo errore genera anche un segnale \signal{SIGPIPE}, la cui azione
  predefinita è terminare il programma; pertanto non si potrà vedere questo
  errore fintanto che \signal{SIGPIPE} non viene gestito o bloccato.
\item[\errcode{EREMOTE} \textit{Object is remote}.] Si è fatto un tentativo di
  montare via NFS un filesystem remoto con un nome  che già specifica un
  filesystem montato via NFS. 
\item[\errcode{EROFS} \textit{Read-only file system}.]  Si è cercato di
  eseguire una operazione di scrittura su un file o una directory che risiede
  su un filesystem montato un sola lettura.
\item[\errcode{ESPIPE} \textit{Invalid seek operation}.] Si cercato di
  eseguire una \func{lseek} su un file che non supporta questa operazione (ad
  esempio su una \textit{pipe}, da cui il nome).
\item[\errcode{ESTALE} \textit{Stale NFS file handle}.] Indica un problema
  interno a NFS causato da cambiamenti del filesystem del sistema remoto. Per
  recuperare questa condizione in genere è necessario smontare e rimontare il
  filesystem NFS.
\item[\errcode{ETXTBSY} \textit{Text file busy}.] Si è cercato di eseguire un
  file che è aperto in scrittura, o di scrivere su un file che è in
  esecuzione.
\item[\errcode{EUSERS} \textit{Too many users}.] Troppi utenti, il sistema delle
  quote rileva troppi utenti nel sistema.
\item[\errcode{EXDEV} \textit{Cross-device link}.] Si è tentato di creare un
  link diretto che attraversa due filesystem differenti.
\end{basedescript}

\section{Gli errori dei processi}
\label{sec:err_proc_errors}

In questa sezione sono raccolti i codici restituiti dalle \textit{system call}
attinenti ad errori che riguardano operazioni specifiche relative alla
gestione dei processi.

\begin{basedescript}{\desclabelwidth{1.5cm}\desclabelstyle{\nextlinelabel}}
\item[\errcode{E2BIG} \textit{Argument list too long}.] La lista degli
  argomenti passati è troppo lunga: è una condizione prevista da POSIX quando
  la lista degli argomenti passata ad una delle funzioni \func{exec} occupa
  troppa memoria.
\item[\errcode{ECHILD} \textit{There are no child processes}.] Non esistono
  processi figli di cui attendere la terminazione. Viene rilevato dalle
  funzioni \func{wait} e \func{waitpid} (vedi sez.~\ref{sec:proc_wait}).
\item[\errcode{EPROCLIM} \textit{Too many processes}.]  Il limite dell'utente
  per nuovi processi (vedi sez.~\ref{sec:sys_resource_limit}) sarà ecceduto
  alla prossima \func{fork}; è un codice di errore di BSD, che non viene
  utilizzato al momento su Linux. 
\item[\errcode{ESRCH} \textit{No process matches the specified process ID}.]
  Non esiste un processo o un \textit{process group} corrispondenti al valore
  dell'identificativo specificato.
\end{basedescript}


\section{Gli errori di rete}
\label{sec:err_network}

In questa sezione sono raccolti i codici restituiti dalle \textit{system call}
attinenti ad errori che riguardano operazioni specifiche relative alla
gestione dei socket e delle connessioni di rete.

\begin{basedescript}{\desclabelwidth{1.5cm}\desclabelstyle{\nextlinelabel}}
\item[\errcode{EADDRINUSE} \textit{Address already in use}.] L'indirizzo del
  socket richiesto è già utilizzato (ad esempio si è eseguita \func{bind}
  su una porta già in uso).
\item[\errcode{EADDRNOTAVAIL} \textit{Cannot assign requested
    address}.]  L'indirizzo richiesto non è disponibile (ad esempio si
  è cercato di dare al socket un nome che non corrisponde al nome
  della stazione locale), o l'interfaccia richiesta non esiste.
\item[\errcode{EAFNOSUPPORT} \textit{Address family not supported by
    protocol}.]  Famiglia di indirizzi non supportata. La famiglia di
  indirizzi richiesta non è supportata, o è inconsistente con il protocollo
  usato dal socket.
\item[\errcode{ECONNABORTED} \textit{Software caused connection abort}.] Una
  connessione è stata abortita localmente. 
\item[\errcode{ECONNREFUSED} \textit{Connection refused}.] Un host remoto ha
  rifiutato la connessione (in genere dipende dal fatto che non c'è un server
  per soddisfare il servizio richiesto).  
\item[\errcode{ECONNRESET} \textit{Connection reset by peer}.] Una connessione
  è stata chiusa per ragioni fuori dal controllo dell'host locale, come il
  riavvio di una macchina remota o un qualche errore non recuperabile sul
  protocollo.
\item[\errcode{EDESTADDRREQ} \textit{Destination address required}.] Non c'è
  un indirizzo di destinazione predefinito per il socket. Si ottiene questo
  errore mandando dato su un socket senza connessione senza averne prima
  specificato una destinazione.
\item[\errcode{EISCONN} \textit{Transport endpoint is already connected}.] Si
  è tentato di connettere un socket che è già connesso.
\item[\errcode{EMSGSIZE} \textit{Message too long}.] Le dimensioni di un
  messaggio inviato su un socket sono eccedono la massima lunghezza supportata.
\item[\errcode{ENETDOWN} \textit{Network is down}.] L'operazione sul socket è
  fallita perché la rete è sconnessa.
\item[\errcode{ENETRESET} \textit{Network dropped connection because of
    reset}.]  Una connessione è stata cancellata perché l'host remoto è
  caduto.
\item[\errcode{ENETUNREACH} \textit{Network is unreachable}.] L'operazione è
  fallita perché l'indirizzo richiesto è irraggiungibile (ad esempio la
  sottorete della stazione remota è irraggiungibile).
\item[\errcode{ENOBUFS} \textit{No buffer space available}.] Tutti i buffer per
  le operazioni di I/O del kernel sono occupati. In generale questo errore è
  sinonimo di \errcode{ENOMEM}, ma attiene alle funzioni di input/output. In
  caso di operazioni sulla rete si può ottenere questo errore invece
  dell'altro.
\item[\errcode{ENOPROTOOPT} \textit{Protocol not available}.] Protocollo non
  disponibile. Si è richiesta un'opzione per il socket non disponibile con il
  protocollo usato.
\item[\errcode{ENOTCONN} \textit{Transport endpoint is not connected}.] Il
  socket non è connesso a niente. Si ottiene questo errore quando si cerca di
  trasmettere dati su un socket senza avere specificato in precedenza la loro
  destinazione. Nel caso di socket senza connessione (ad esempio socket UDP)
  l'errore che si ottiene è \errcode{EDESTADDRREQ}.
\item[\errcode{ENOTSOCK} \textit{Socket operation on non-socket}.] Si è
  tentata un'operazione su un file descriptor che non è un socket quando
  invece era richiesto un socket.
\item[\errcode{EOPNOTSUPP} \textit{Operation not supported on transport
    endpoint}.] L'operazione richiesta non è supportata. Alcune funzioni non
  hanno senso per tutti i tipi di socket, ed altre non sono implementate per
  tutti i protocolli di trasmissione. Questo errore quando un socket non
  supporta una particolare operazione, e costituisce una indicazione generica
  che il server non sa cosa fare per la chiamata effettuata.
\item[\errcode{EPFNOSUPPORT} \textit{Protocol family not supported}.] Famiglia
  di protocolli non supportata. La famiglia di protocolli richiesta non è
  supportata.
\item[\errcode{EPROTONOSUPPORT} \textit{Protocol not supported}.] Protocollo
  non supportato. Il tipo di socket non supporta il protocollo richiesto (un
  probabile errore nella specificazione del protocollo).
\item[\errcode{EPROTOTYPE} \textit{Protocol wrong type for socket}.]
  Protocollo sbagliato per il socket. Il socket usato non supporta il
  protocollo di comunicazione richiesto.
\item[\errcode{ESHUTDOWN} \textit{Cannot send after transport endpoint
    shutdown}.] Il socket su cui si cerca di inviare dei dati ha avuto uno
  shutdown.
\item[\errcode{ESOCKTNOSUPPORT} \textit{Socket type not supported}.] Socket
  non supportato. Il tipo di socket scelto non è supportato.
\item[\errcode{ETIMEDOUT} \textit{Connection timed out}.] Un'operazione sul
  socket non ha avuto risposta entro il periodo di timeout.
\item[\errcode{ETOOMANYREFS} \textit{Too many references: cannot splice}.] La
  \acr{glibc} dice ???
\end{basedescript}


\section{Errori generici}

In questa sezione sono raccolti i codici restituiti dalle \textit{system call}
attinenti ad errori generici, si trovano qui tutti i codici di errore non
specificati nelle sezioni precedenti.

\begin{basedescript}{\desclabelwidth{1.5cm}\desclabelstyle{\nextlinelabel}}
\item[\errcode{EAGAIN} \textit{Resource temporarily unavailable}.] La funzione è
  fallita ma potrebbe funzionare se la chiamata fosse ripetuta. Questo errore
  accade in due tipologie di situazioni:
  \begin{itemize}
  \item Si è effettuata un'operazione che si sarebbe bloccata su un oggetto
    che è stato posto in modalità non bloccante. Nei vecchi sistemi questo era
    un codice diverso, \errcode{EWOULDBLOCK}. In genere questo ha a che fare
    con file o socket, per i quali si può usare la funzione \func{select} per
    vedere quando l'operazione richiesta (lettura, scrittura o connessione)
    diventa possibile.
  \item Indica la carenza di una risorsa di sistema che non è al momento
    disponibile (ad esempio \func{fork} può fallire con questo errore se si è
    esaurito il numero di processi contemporanei disponibili). La ripetizione
    della chiamata in un periodo successivo, in cui la carenza della risorsa
    richiesta può essersi attenuata, può avere successo. Questo tipo di
    carenza è spesso indice di qualcosa che non va nel sistema, è pertanto
    opportuno segnalare esplicitamente questo tipo di errori.
  \end{itemize}
\item[\errcode{EALREADY} \textit{Operation already in progress}.] L'operazione è
  già in corso. Si è tentata un'operazione già in corso su un oggetto posto in
  modalità non-bloccante.
\item[\errcode{EDEADLK} \textit{Deadlock avoided}.] L'allocazione di una
  risorsa avrebbe causato un \textit{deadlock}. Non sempre il sistema è in
  grado di riconoscere queste situazioni, nel qual caso si avrebbe il blocco.
\item[\errcode{EFAULT} \textit{Bad address}.] Una stringa passata come
  argomento è fuori dello spazio di indirizzi del processo, in genere questa
  situazione provoca direttamente l'emissione di un segnale di \textit{segment
    violation} (\signal{SIGSEGV}).
\item[\errcode{EDOM} \textit{Domain error}.] È usato dalle funzioni matematiche
  quando il valore di un argomento è al di fuori dell'intervallo in cui esse
  sono definite.
\item[\errcode{EILSEQ} \textit{Illegal byte sequence}.] Nella decodifica di un
  carattere esteso si è avuta una sequenza errata o incompleta o si è
  specificato un valore non valido.
\item[\errcode{EINPROGRESS} \textit{Operation now in progress}.] Operazione in
  corso. Un'operazione che non può essere completata immediatamente è stata
  avviata su un oggetto posto in modalità non-bloccante. Questo errore viene
  riportato per operazioni che si dovrebbero sempre bloccare (come per una
  \func{connect}) e che pertanto non possono riportare \errcode{EAGAIN},
  l'errore indica che l'operazione è stata avviata correttamente e occorrerà
  del tempo perché si possa completare. La ripetizione della chiamata darebbe
  luogo ad un errore \errcode{EALREADY}.
\item[\errcode{EINTR} \textit{Interrupted function call}.] Una funzione di
  libreria è stata interrotta. In genere questo avviene causa di un segnale
  asincrono al processo che impedisce la conclusione della chiamata, la
  funzione ritorna con questo errore una volta che si sia correttamente
  eseguito il gestore del segnale. In questo caso è necessario ripetere la
  chiamata alla funzione.
\item[\errcode{EINVAL} \textit{Invalid argument}.] Errore utilizzato per
  segnalare vari tipi di problemi dovuti all'aver passato un argomento
  sbagliato ad una funzione di libreria.
\item[\errcode{ENOMEM} \textit{No memory available}.] Il kernel non è in grado
  di allocare ulteriore memoria per completare l'operazione richiesta.
\item[\errcode{ENOSYS} \textit{Function not implemented}.] Indica che la
  funzione non è supportata o nelle librerie del C o nel kernel. Può dipendere
  sia dalla mancanza di una implementazione, che dal fatto che non si è
  abilitato l'opportuno supporto nel kernel; nel caso di Linux questo può
  voler dire anche che un modulo necessario non è stato caricato nel sistema.
\item[\errcode{ENOTSUP} \textit{Not supported}.] Una funzione ritorna questo
  errore quando gli argomenti sono validi ma l'operazione richiesta non è
  supportata. Questo significa che la funzione non implementa quel particolare
  comando o opzione o che, in caso di oggetti specifici (file descriptor o
  altro) non è in grado di supportare i parametri richiesti.
\item[\errcode{ERANGE} \textit{Range error}.] È usato dalle funzioni
  matematiche quando il risultato dell'operazione non è rappresentabile nel
  valore di ritorno a causa di un overflow o di un underflow.
\item[\errcode{EWOULDBLOCK} \textit{Operation would block}.] Indica che
  l'operazione richiesta si bloccherebbe, ad esempio se si apre un file in
  modalità non bloccante, una \func{read} restituirebbe questo errore per
  indicare che non ci sono dati; in Linux è identico a \errcode{EAGAIN}, ma in
  altri sistemi può essere specificato un valore diverso.
\end{basedescript}


\begin{basedescript}{\desclabelwidth{1.5cm}\desclabelstyle{\nextlinelabel}}
% definiti nel manuale delle glibc ma inesistenti in linux/errno.h
%\item[\errcode{EBADRPC} \textit{}.] 
%\item[\errcode{ERPCMISMATCH} \textit{}.] 
%\item[\errcode{EPROGUNAVAIL} \textit{}.] 
%\item[\errcode{EPROGMISMATCH} \textit{}.] 
%\item[\errcode{EPROCUNAVAIL} \textit{}.] 
%\item[\errcode{EAUTH} \textit{}.] 
%\item[\errcode{ENEEDAUTH} \textit{}.] 
%\item[\errcode{EBACKGROUND} \textit{}.] 
%\item[\errcode{EDIED} \textit{}.] 
% questi sembrano scherzi, sempre dal manuale delle glibc...
%\item[\errcode{ED} \textit{}.] 
%\item[\errcode{EGREGIOUS} \textit{}.] 
%\item[\errcode{EIEIO} \textit{}.] 
%\item[\errcode{EGRATUITOUS} \textit{} roba di Hurd, pare. 


\item[\errcode{EBADMSG} \textit{Not a data message}.] Definito da POSIX come
errore che arriva ad una funzione di lettura che opera su uno stream. Non
essendo gli stream definiti su Linux il kernel non genera mai questo tipo di
messaggio. 

\item[\errcode{EIDRM} \textit{Identifier removed}.] Indica che l'oggetto del
  \textit{SysV IPC} a cui si fa riferimento è stato cancellato.

\item[\errcode{EMULTIHOP} \textit{Multihop attempted}.] Definito da POSIX come
  errore dovuto all'accesso a file remoti attraverso più macchine, quando ciò
  non è consentito. Non viene mai generato su Linux.

\item[\errcode{ENOATTR} \textit{No such attribute}.] È un codice di errore
  specifico di Linux utilizzato dalle funzioni per la gestione degli attributi
  estesi dei file (vedi sez.~\ref{sec:file_xattr}) quando il nome
  dell'attributo richiesto non viene trovato.

\item[\errcode{ENODATA} \textit{No data available}.] Viene indicato da POSIX
  come restituito da una \func{read} eseguita su un file descriptor in
  modalità non bloccante quando non ci sono dati. In realtà in questo caso su
  Linux viene utilizzato \errcode{EAGAIN}. Lo stesso valore però viene usato
  come sinonimo di \errcode{ENOATTR}.

\item[\errcode{ENOLINK} \textit{Link has been severed}.] È un errore il cui
  valore è indicato come \textsl{riservato} nelle \textit{Single Unix
    Specification}. Dovrebbe indicare l'impossibilità di accedere ad un file a
  causa di un errore sul collegamento di rete, ma non ci sono indicazioni
  precise del suo utilizzo. Per quanto riguarda Linux viene riportato nei
  sorgenti del kernel in alcune operazioni relative ad operazioni di rete. 

\item[\errcode{ENOMSG} \textit{No message of desired type}.] Indica che in una
  coda di messaggi del \textit{SysV IPC} non è presente nessun messaggio del
  tipo desiderato.

\item[\errcode{ENOSR} \textit{Out of streams resources}.] Errore relativo agli
  \textit{STREAMS}, che indica l'assenza di risorse sufficienti a completare
  l'operazione richiesta. Quella degli \textit{STREAMS}\footnote{che non vanno
    confusi con gli \textit{stream} di sez.~\ref{sec:files_std_interface}.}  è
  interfaccia di programmazione originaria di System V, che non è implementata
  da Linux, per cui questo errore non viene utilizzato.

\item[\errcode{ENOSTR} \textit{Device not a stream}.] Altro errore relativo
  agli \textit{STREAMS}, anch'esso non utilizzato da Linux.

\item[\errcode{EOVERFLOW} \textit{Value too large for defined data type}.] Si è
  chiesta la lettura di un dato dal \textit{SysV IPC} con \const{IPC\_STAT} ma
  il valore eccede la dimensione usata nel buffer di lettura.

\item[\errcode{EPROTO} \textit{Protocol error}.] Indica che c'è stato un errore
  nel protocollo di rete usato dal socket; viene usato come errore generico
  dall'interfaccia degli \textit{STREAMS} quando non si è in grado di
  specificare un altro codice di errore che esprima più accuratamente la
  situazione.

\item[\errcode{ETIME} \textit{Timer expired}.] Indica che è avvenuto un timeout
  nell'accesso ad una risorsa (ad esempio un semaforo). Compare nei sorgenti
  del kernel (in particolare per le funzioni relativa al bus USB) come
  indicazione di una mancata risposta di un dispositivo, con una descrizione
  alternativa di \textit{Device did not respond}.
\end{basedescript}



% \section{Errori del kernel}
% \label{sec:err_kernel_err}

% In questa sezione sono raccolti i codici di errore interni del kernel. Non
% sono usati dalle funzioni di libreria, ma vengono riportati da alcune system
% call 
% TODO verificare i dettagli degli errori del kernel, eventualmente cassare.

% \begin{description}
% \item[\errcode{ERESTART} \textit{Interrupted system call should be restarted}.] 
% \item[\errcode{ECHRNG} \textit{Channel number out of range}.] 
% \item[\errcode{EL2NSYNC} \textit{Level 2 not synchronized}.] 
% \item[\errcode{EL3HLT} \textit{Level 3 halted}.] 
% \item[\errcode{EL3RST} \textit{Level 3 reset}.] 
% \item[\errcode{ELNRNG} \textit{Link number out of range}.] 
% \item[\errcode{EUNATCH} \textit{Protocol driver not attached}.] 
% \item[\errcode{ENOCSI} \textit{No CSI structure available}.] 
% \item[\errcode{EL2HLT} \textit{Level 2 halted}.] 
% \item[\errcode{EBADE} \textit{Invalid exchange}.] 
% \item[\errcode{EBADR} \textit{Invalid request descriptor}.] 
% \item[\errcode{EXFULL} \textit{Exchange full}.] 
% \item[\errcode{ENOANO} \textit{No anode}.] 
% \item[\errcode{EBADRQC} \textit{Invalid request code}.] 
% \item[\errcode{EBADSLT} \textit{Invalid slot}.] 
% \item[\errcode{EDEADLOCK} Identico a \errcode{EDEADLK}.] 
% \item[\errcode{EBFONT} \textit{Bad font file format}.] 
% \item[\errcode{ENONET} \textit{Machine is not on the network}.] 
% \item[\errcode{ENOPKG} \textit{Package not installed}.] 
% \item[\errcode{EADV} \textit{Advertise error}.] 
% \item[\errcode{ESRMNT} \textit{Srmount error}.] 
% \item[\errcode{ECOMM} \textit{Communication error on send}.] 
% \item[\errcode{EDOTDOT} \textit{RFS specific error}.] 
% \item[\errcode{ENOTUNIQ} \textit{Name not unique on network}.] 
% \item[\errcode{EBADFD} \textit{File descriptor in bad state}.] 
% \item[\errcode{EREMCHG} \textit{Remote address changed}.] 
% \item[\errcode{ELIBACC} \textit{Can not access a needed shared library}.] 
% \item[\errcode{ELIBBAD} \textit{Accessing a corrupted shared library}.] 
% \item[\errcode{ELIBSCN} \textit{.lib section in a.out corrupted}.] 
% \item[\errcode{ELIBMAX} \textit{Attempting to link in too many shared
%     libraries}.]
% \item[\errcode{ELIBEXEC} \textit{Cannot exec a shared library directly}.] 
% \item[\errcode{ESTRPIPE} \textit{Streams pipe error}.] 
% \item[\errcode{EUCLEAN} \textit{Structure needs cleaning}.] 
% \item[\errcode{ENAVAIL} \textit{No XENIX semaphores available}.] 
% \item[\errcode{EISNAM} \textit{Is a named type file}.] 
% \item[\errcode{EREMOTEIO} \textit{Remote I/O error}.] 
% \item[\errcode{ENOMEDIUM} \textit{No medium found}.] 
% \item[\errcode{EMEDIUMTYPE} \textit{Wrong medium type}.] 
% \end{description}


% LocalWords:  header errno perror int strerror sez EPERM Operation not ENOENT
% LocalWords:  permitted such pathname EIO error ENXIO device address kernel Is
% LocalWords:  ENOEXEC Invalid executable format exec EBADF Bad descriptor Too
% LocalWords:  EACCES Permission denied ELOOP many symbolic links encountered
% LocalWords:  ENAMETOOLONG name too long ENOTBLK Block required block EEXIST
% LocalWords:  exists EBUSY Resource busy filesystem EXDEV ENODEV ENOTDIR files
% LocalWords:  EISDIR EMFILE ulimit ENFILE table overflow ENOTTY ETXTBSY Text
% LocalWords:  EFBIG big ENOSPC left ESPIPE seek operation EROFS Read only read
% LocalWords:  system EMLINK EPIPE Broken SIGPIPE ENOTEMPTY empty EUSERS users
% LocalWords:  EDQUOT exceeded ESTALE NFS handle EREMOTE Object is ENOLCK locks
% LocalWords:  available locking EFTYPE type sticky ESRCH process matches the
% LocalWords:  specified pid Argument list POSIX ECHILD There child processes
% LocalWords:  socket ENOTSOCK EMSGSIZE Message EPROTOTYPE Protocol wrong for
% LocalWords:  ENOPROTOOPT EPROTONOSUPPORT supported ESOCKTNOSUPPORT EOPNOTSUPP
% LocalWords:  transport endpoint EPFNOSUPPORT family EAFNOSUPPORT protocol of
% LocalWords:  EADDRINUSE already bind EADDRNOTAVAIL Cannot assign requested to
% LocalWords:  ENETDOWN ENETUNREACH unreachable ENETRESET dropped connection
% LocalWords:  because reset l'host ECONNABORTED caused abort ECONNRESET peer
% LocalWords:  dell'host ENOBUFS ENOMEM EISCONN connected ENOTCONN UDP send now
% LocalWords:  EDESTADDRREQ Destination ESHUTDOWN after shutdown ETOOMANYREFS
% LocalWords:  references cannot splice glibc ETIMEDOUT timed ECONNREFUSED host
% LocalWords:  refused EHOSTDOWN EHOSTUNREACH route EINTR Interrupted function
% LocalWords:  call memory EDEADLK Deadlock avoided deadlock EFAULT segment IPC
% LocalWords:  violation SIGSEGV EINVAL argument EDOM Domain ERANGE underflow
% LocalWords:  EAGAIN temporarily unavailable EWOULDBLOCK select fork would has
% LocalWords:  EINPROGRESS progress connect EALREADY ENOSYS implemented ENOTSUP
% LocalWords:  EILSEQ Illegal sequence EBADMSG message EIDRM Identifier removed
% LocalWords:  SysV EMULTIHOP Multihop attempted ENODATA ENOLINK been severed
% LocalWords:  ENOMSG desired ENOSR streams resources ENOSTR stream EOVERFLOW
% LocalWords:  Value large defined STAT EPROTO ETIME Timer expired group wait
% LocalWords:  waitpid Specification cap USB did respond Stale


%%% Local Variables: 
%%% mode: latex
%%% TeX-master: "gapil"
%%% End: 
